%%%%%%%%%%  scan idx 129  |  folio 116  %%%%%%%%%%
\SGAPrintPage{116}%
\uline{In particular, for} \(F\) \uline{to be pseudo-coherent, it is necessary and sufficient that} \(\Ho^n(F)\) \uline{be coherent for every} \(n\)\uline{.}

\uline{Proof.} Assertion a) and the implication \(\Rightarrow\) of b) follow trivially from (3.4). The implication \(\Leftarrow\) of b) follows at once from a) and (2.11.b)).

\sgaStableTarget{sga6.e01.ex.3-6}\textbf{Examples 3.6.} Let \((\underline{S},A)\) be a ringed topos, and \(\underline{C},\underline{C}_0\) the \(\underline{S}\)-categories of (1.3.d)) (cf.\ also (2.15)). For the objects of \(\underline{C}_0\) to be coherent, it is necessary and sufficient that \(A\) be coherent (i.e.\ precoherent). Here are examples where this holds:
\begin{itemize}[label=-]
\item \(A\) is constant with value a left noetherian ring (SGA 4 IX 2.1).
\item \(\underline{S}\) is the topos of sheaves of sets (for the Zariski topology) on a locally noetherian scheme \(S\)\SGAAuthorFootnote{*}{In accordance with the new conventions, we shall say ``scheme'' for ``prescheme'', and ``separated scheme'' for ``scheme''.}, and \(A=\underline{O}_S\).
\item \(\underline{S}\) is the topos of sheaves of sets on a complex analytic space \(S\), \(A=\underline{O}_S\).
\end{itemize}
On the other hand, the sheaf of continuous complex-valued functions on a topological space is in general not coherent.

Of course, when \((\underline{S},A)\) is defined by a ringed space, the notion of coherence considered here coincides with Serre's usual notion (EGA \(0_I\) 5.3.1).

\sgaStableTarget{sga6.e01.prop.3-7}\textbf{Proposition 3.7.} a) \uline{Let} \(f:(\underline{S}',A')\to(\underline{S},A)\) \uline{be a morphism of ringed topoi. Assume that} \(f\) \uline{is flat on the right,} i.e.\ \uline{that the functor} \(f^*=A'\otimes_A.\) \uline{is exact. For} \(F\in\Ob\Dpx(_A S)\)\uline{,} \(G\in\Ob\Dpx^+(_A S)\)\uline{, there is a functorial canonical arrow:}
\[
  f^*\mathbb{R}\underline{\Hom}(F,G)\longrightarrow\mathbb{R}\underline{\Hom}(f^*(F),f^*(G)).
\]
%%%%%%%%%%  scan idx 130  |  folio 117  %%%%%%%%%%
\SGAPrintPage{117}%
\uline{This is an isomorphism when} \(F\) \uline{is pseudo-coherent.}

b) \uline{Suppose that} \(A\) \uline{is commutative and coherent, and let} \(F\in\Ob\Dpx^-(S)\)\uline{,} \(G\in\Ob\Dpx^+(S)\)\uline{. If} \(\underline{\Ho}^i(F)\) and \(\underline{\Ho}^i(G)\) \uline{are coherent for every} \(i\)\uline{, the same is true of} \(\underline{\Ext}^i(F,G)\)\uline{.}

\uline{Proof.} a) To define the arrow (cf.\ SGA 4 XVII), we can suppose that \(G\) has injective components, and take the composite arrow
\[
  f^*\underline{\Hom}(F,G)\longrightarrow\underline{\Hom}(f^*F,f^*G)\longrightarrow\mathbb{R}\underline{\Hom}(f^*F,f^*G).
\]
Let us show that it is an isomorphism for \(F\) pseudo-coherent. For brevity, write \(T\) (resp.\ \(T'\)) for the functor \(f^*\mathbb{R}\underline{\Hom}(.,G)\) (resp.\ \(\mathbb{R}\underline{\Hom}(f^*.,f^*G)\)), and let \(a\in\mathbb Z\) be such that \(\Ho^i(G)=0\) for \(i<a\). It is clear that \(TA\qis T'A\), hence, by dévissage, \(TF\qis T'F\) for \(F\) strictly perfect (2.1). On the other hand, when \(F\) is acyclic in degree \(>n\), \(TF\) and \(T'F\) are acyclic in degree \(<a-n\). To verify that \(TF\qis T'F\) for \(F\) pseudo-coherent, we can, since the question is local, suppose that \(F\) is strictly \(n\)-pseudo-coherent, i.e.\ that there is a distinguished triangle
\[
\begin{tikzcd}[column sep=small]
& F'' \arrow[dl,"+1"'] & \\
F' \arrow[rr] & & F \arrow[ul]
\end{tikzcd}
\;,
\]
with \(F'\) strictly perfect and \(F''\) acyclic in degree \(>n-1\). We deduce a morphism of distinguished triangles
\[
\begin{tikzcd}[column sep=small]
& TF'' \arrow[dr] & \\
TF' \arrow[ur,"-1"'] & & TF \arrow[ll]
\end{tikzcd}
\qquad\longrightarrow\qquad
\begin{tikzcd}[column sep=small]
& T'F'' \arrow[dr] & \\
T'F' \arrow[ur,"-1"'] & & T'F \arrow[ll]
\end{tikzcd}
\;.
\]
Now \(TF'\qis T'F'\) since \(F'\) is strictly perfect, and \(TF''\) and \(T'F''\) are
%%%%%%%%%%  scan idx 131  |  folio 118  %%%%%%%%%%
\SGAPrintPage{118}%
acyclic in degree \(\leq a-n\); writing the induced morphism on the exact cohomology sequences of the preceding triangles, we deduce that \(TF\to T'F\) induces an isomorphism \(\Ho^i(TF)\qis\Ho^i(T'F)\) for \(i\leq a-n\). Since we can choose \(n\) arbitrarily negative, we are done.

b) The proof is left as an exercise to the reader: see, for example, ([H] II 3.3).

\sgaStableTarget{sga6.e01.rem.3-7-1}\textbf{Remark 3.7.1.} We shall see below (7.1.2) another case of ``compatibility'' of \(\mathbb{R}\underline{\Hom}\) with \(\mathbb{L}f^*\).

\sgaStableTarget{sga6.e01.cor.3-7-2}\textbf{Corollary 3.7.2.} (cf.\ EGA \(0_I\) 5.2.7). \uline{Let} \((\underline{S},A)\) \uline{be a ringed topos,} \(s\) \uline{a point of} \(\underline{S}\) (SGA 4 VIII 7.8)\uline{,} \(E,F\in\Ob\D^b(_A S)_{\coh}\)\uline{. If} \(E_s\simeq F_s\)\uline{, there exists an} \(s\)\uline{-pointed object} \(U\) \uline{of} \(\underline{S}\) \uline{such that} \(E|U\simeq F|U\)\uline{.}

\uline{Proof.} Let \(\alpha:E_s\to F_s\) and \(\beta:F_s\to E_s\) be reciprocal isomorphisms. By (3.7 a)), we have
\[
  \underline{\Hom}(E,F)_s\qis\Hom(E_s,F_s)
  \qquad
  (\text{resp. }\underline{\Hom}(F,E)_s\qis\Hom(F_s,E_s)).
\]
There are therefore a ``neighbourhood'' \(U\) of \(s\) and a section \(u\) (resp.\ \(v\)) of \(\underline{\Hom}(E,F)\) (resp.\ \(\underline{\Hom}(F,E)\)) over \(U\) such that \(u_s=\alpha\) (resp.\ \(v_s=\beta\)). The homomorphisms \(vu:E|U\to E|U\) and \(uv:F|U\to F|U\) induce the identity at \(s\), hence over a suitable \(V\) over \(U\), as required.

%%%%%%%%%%  scan idx 132  |  folio 119  %%%%%%%%%%
\SGAPrintPage{119}%
\sgaStableTarget{sga6.e01.sec.4}\section*{4. Perfect complexes}

\sgaStableTarget{sga6.e01.para.4-0}4.0. In this section, \(\underline{S}\) denotes a site, \(S\) an object of \(\underline{S}\), \(\underline{C}\) a flat abelian \(\underline{S}\)-category, and \(\underline{C}_0\) a strictly full additive sub-\(\underline{S}\)-category of \(\underline{C}\) (1.1). Unless otherwise mentioned, we assume that \(\underline{C}_0\) is locally liftable in \(\underline{C}\) and locally stable under kernels of epimorphisms (1.2).

\sgaStableTarget{sga6.e01.lem.4-1}\textbf{Lemma 4.1}\SGAAuthorFootnote{raisedstar}{This lemma was communicated to the editor by D. Ferrand.}. \uline{Let} \(u:E\to F\) \uline{be an arrow of} \(\Cpx(\underline{C}_S)\)\uline{. If} \(E\) \uline{is strictly} \(\underline{C}_0\)\uline{-perfect} (2.1) \uline{and} \(F\) \uline{is acyclic, then} \(u\) \uline{is locally homotopic to zero.}

\uline{Proof.} We must construct locally a homotopy \(k\) such that \(u=dk+kd\). Since \(E\) is strictly perfect, there exist \(a,b\in\mathbb Z\), with \(a\leq b\), such that \(E^i=0\) for \(i\notin[a,b]\). We can therefore take \(k^i=0\) for \(i\notin[a,b]\). We then construct \(k^i\) for \(i\geq a\) by descending induction on \(i\). Suppose that \(k^i:E^i\to F^{i-1}\) has been constructed so that \(u^i=k^{i+1}d_E^i+d_F^{i-1}k^i\) for \(i\geq n>a\), and let us show that, after localizing, there exists \(k^{n-1}:E^{n-1}\to F^{n-2}\) such that \(u^{n-1}=k^nd_E^{n-1}+d_F^{n-2}k^{n-1}\). By the induction hypothesis, the arrow
\[
  v^{n-1}=u^{n-1}-k^nd_E^{n-1}:E^{n-1}\to F^{n-1}
\]
factors through \(Z^{n-1}(F)\). But \(Z^{n-1}(F)=B^{n-1}(F)\) since \(F\) is acyclic. Since \(\underline{C}_0\) is locally liftable in \(\underline{C}\), it follows that, after localizing, there exists \(k^{n-1}:E^{n-1}\to F^{n-2}\) such that \(v^{n-1}=d_F^{n-2}k^{n-1}\), as required.

\sgaStableTarget{sga6.e01.cor.4-2}\textbf{Corollary 4.2.} \uline{Let}
\[
\begin{tikzcd}
& P \arrow[d,"v"] \\
E \arrow[r,"f"] & F
\end{tikzcd}
\]
%%%%%%%%%%  scan idx 133  |  folio 120  %%%%%%%%%%
\SGAPrintPage{120}%
\uline{be a diagram of} \(\Cpx(\underline{C}_S)\)\uline{, where} \(P\) \uline{is strictly} \(\underline{C}_0\)\uline{-perfect and} \(f\) \uline{a quasi-isomorphism. The set of} \(X\) \uline{over} \(S\) \uline{for which there exists an arrow} \(u:P_X\to E_X\) \uline{of} \(\Cpx(\underline{C}_X)\) \uline{such that the diagram}
\[
\begin{tikzcd}
& P_X \arrow[d,"v_X"] \arrow[dl,"u"']\\
E_X \arrow[r,"f_X"] & F_X
\end{tikzcd}
\]
\uline{is commutative in} \(\Kpx(\underline{C}_X)\) \uline{is a refinement of} \(S\)\uline{.}

\uline{Proof.} Let \(G\) be the cone of \(f\). We have the exact sequence
\[
  \longrightarrow \Hom_{\Kpx(\underline{C}_S)}(P,E)
  \longrightarrow \Hom_{\Kpx(\underline{C}_S)}(P,F)
  \longrightarrow \Hom_{\Kpx(\underline{C}_S)}(P,G)
  \longrightarrow .
\]
Since \(G\) is acyclic, the image of \(v\) in \(\Hom_{\Kpx(\underline{C}_S)}(P,G)\) is locally zero by (4.1). Thus, after localizing, there exists \(u:P\to E\) such that \(v=fu\) in \(\Kpx(\underline{C}_S)\), as required.

\sgaStableTarget{sga6.e01.cor.4-3}\textbf{Corollary 4.3.} \uline{Let} \(f:E\to F\) \uline{be a quasi-isomorphism of} \(\Cpx(\underline{C}_S)\)\uline{, with} \(F\) \uline{strictly perfect (relative to} \(\underline{C}_0\)\uline{). The set of} \(X\) \uline{over} \(S\) \uline{for which there exists a quasi-isomorphism} \(g:F_X\to E_X\) \uline{such that} \(f_Xg\) \uline{is homotopic to} \(\Id_{F_X}\) \uline{is a refinement of} \(S\)\uline{.}

\uline{Proof.} This is a special case of (4.2).

\sgaStableTarget{sga6.e01.schol.4-4}\textbf{Scholium 4.4.} Let \(F\in\Ob\Kpx(\underline{C}_S)\). We can define (for example using a cleavage of \(\underline{C}\) (SGA 1 VI 7.1)) a category \((\underline{\Qis/F})\) fibered over \(\underline{S}_{/S}\), whose fibre at \(X\) over \(S\) is the category \((\Qis/F_X)^\circ\) (opposite to the category of quasi-isomorphisms with target \(F_X\)). We can then express (4.3) by saying that,
